% ============================================================================
% I. INTRODUCTION  (target length: about 0.75 column)
% Source keys: see sections/00_abstract.tex.
% Positioning after the verified novelty search (2026-10-16): early and
% speculative retrieval on partial transcripts is PRIOR WORK (Stream RAG,
% speculative execution for voice agents); the paper builds on it. Novelty
% claims are limited to (i) the formal joint-vs-decomposed analysis and the
% pre-registered boundary test, (ii) the persona finding on partial
% transcripts, (iii) the matched-pair spoken benchmark -- each hedged "to our
% knowledge". Never write "Speculative RAG"; write "speculative prefetch".
% ============================================================================
\section{Introduction}\label{sec:intro}

% ---- P-I.1 ------------------------------------------------------------------
% PURPOSE: the systems problem. A voice assistant that grounds answers in
%   documents must decide (i) whether to retrieve for this turn and (ii) when
%   to dispatch the retrieval relative to the user's ongoing speech.
% CLAIMS + EVIDENCE:
%   (a) Recent spoken-dialogue systems already retrieve while the user is
%       still speaking, either speculatively at fixed intervals or when a
%       model triggers a call, and report lower latency at preserved accuracy
%       [arora2026streamrag]; speculative execution hides tool latency in
%       cascaded voice agents [jung2026hiding]. This is the setting the paper
%       builds on, not a contribution.
%   (b) Retrieval is not always beneficial: adaptive methods decide whether
%       to retrieve [jiang2023flare, su2024dragin, jeong2024adaptiverag], and
%       retrieval can hurt on questions the model already answers
%       [mallen2023trust].
%   (c) Timing work in this area treats dispatch as safe once the user's
%       intent stabilises and assumes the query needs retrieval
%       [galbraith2026streaming]; whether to retrieve and when are therefore
%       studied separately.
% NUMBERS: none.
\todo{P-I.1: the retrieve-and-when problem in streaming voice RAG, building on
streaming retrieval \cite{arora2026streamrag,jung2026hiding,galbraith2026streaming}
and adaptive retrieval \cite{jiang2023flare,su2024dragin,jeong2024adaptiverag,mallen2023trust}.}

% ---- P-I.2 ------------------------------------------------------------------
% PURPOSE: the question and why it is not obvious. Should the two decisions
%   be taken jointly (one argmax over (retrieve, timing)) or by separate
%   components that share prices? State that the submitted version found a
%   tie and that this revision explains it and identifies when it breaks.
% CLAIMS + EVIDENCE:
%   (a) The pre-registered E3 test found no dominance in 0 of 5 seeds
%       [analysis/e3_dominance.json: seeds_with_strict_dominance = 0].
%   (b) Do NOT present the tie as evidence for or against coupling; present
%       it as the case the theory predicts (A1--A3 hold in v1 by design).
% NUMBERS: 0/5 seeds [analysis/e3_dominance.json].
\todo{P-I.2: the joint-versus-decomposed question; the v1 tie (0/5 seeds)
and the plan to explain it and test its boundaries.}

% ---- P-I.3 ------------------------------------------------------------------
% PURPOSE: contributions, as a numbered list (authors may keep the list
%   form; wording is theirs).
% BACKGROUND, NOT CONTRIBUTIONS (state as adopted, with citations):
%   - uncertainty-triggered retrieval [jiang2023flare, su2024dragin]
%     (R1 minor: "present as adopted");
%   - early and speculative retrieval on partial transcripts, commit-style
%     (model-triggered) versus speculative (fixed-interval) dispatch
%     [arora2026streamrag, jung2026hiding];
%   - self-knowledge probes of the P(IK) kind [kadavath2022language,
%     moskvoretskii2025adaptive].
% CONTRIBUTIONS TO LIST (each hedged "to our knowledge" where novelty is
%   claimed; sources docs/research_v2/NOVELTY.md and THEORY.md as merged at
%   1d8900c, revised after the novelty search):
%   C1  A formal analysis of joint versus decomposed control of retrieval and
%       dispatch time: an equivalence theorem under A1--A3 with an exact price
%       map (Prop. 1) and the fixed-price gap (Prop. 2), which explain the
%       v1 tie; and the conditions under which the equivalence breaks --
%       second order for turn-dependent overlap (Prop. 3), first order for
%       information versus overlap under a COMMIT architecture (Prop. 4), and,
%       under SPECULATIVE prefetch, an advantage that a belief-gated
%       decomposed prefetcher (SPEC-theta) captures up to a C1-sized term
%       (Props. 4', 4''). [TH Secs. 2--6]
%   C2  A pre-registered boundary test of these predictions on the same items
%       under both architectures, whose prediction FAILED (phrase C2 as
%       such): P1 joint > B9+ under commit; P2 the joint advantage is larger
%       under commit than under speculation; P3 under speculation the joint
%       controller is equivalent to the dev-best decomposed rival (B9+ on
%       d1_open, SPEC/B10 on d2_support; candidates B9+, SPEC, SPEC-theta)
%       within +-m3, m3 = 0.016, a smallest effect size of interest fixed
%       before the freeze (TOST); boundary claim iff all three
%       (intersection-union test); plus H-REV, Delta_C < 0, pre-registered
%       before the freeze [PRE2 Secs. 3.2, 3.3]. Outcome on d1_open (the
%       pre-registered per-dataset analysis, unaffected by D3): P1 fails,
%       P2 fails Holm, P3 fails (Delta_S -0.0415, below -m3), H-REV holds
%       (Delta_C -0.0367 [-0.0481, -0.0252], p < 1e-4, survives Holm); the
%       boundary claim is false. d2_support and pooled: exploratory, same
%       direction [docs/research_v2/v2_test_results_final.md].
%   C3  A corrected objective that charges retrieval's measured harm on
%       questions that do not need it (g^- = 0.592 vs g^+ = 0.720 on v1)
%       [V1T retrieval_effect]; together with C1 it explains the sign of the
%       v1 mid-grid gaps.
%   C4  An empirical finding about the belief on partial transcripts in a
%       persona/company-support setting (v2 E1 test result; d2_support is
%       exploratory): entropy-signal failure and P(IK)-style probe
%       sharpening with transcript prefix; evidence (TEST split, frozen analyze.py output
%       llm/experiment1_v2/results/<ds>/analysis_test/analysis_test.json;
%       collected in docs/research_v2/paper_numbers_v2.json e1_test_signals,
%       e1_test_extra): d2_support (EXPLORATORY, judge-graded need label):
%       mean-entropy AUC 0.508 [0.439, 0.576] and max-entropy 0.492
%       [0.419, 0.561], i.e. chance, vs self-assessment probe 0.701
%       [0.638, 0.761]; probe AUC rises with the fraction heard, text
%       0.579/0.607/0.667/0.701 and ASR partials 0.558/0.612/0.672/0.708 at
%       0.3/0.5/0.7/1.0. d1_open (pre-registered, alias-graded): the entropy
%       signal does NOT fail there -- mean-entropy AUC 0.707 [0.691, 0.724]
%       is higher than the probe's 0.665 [0.648, 0.682]; probe text curve
%       0.537/0.543/0.582/0.665 (AUC(1.0) - AUC(0.3) = +0.128
%       [+0.100, +0.155], p < 1e-4; docs/research_v2/v2_test_results_d1.md),
%       ASR curve 0.540/0.561/0.583/0.597. So state the persona finding as
%       specific to the support setting (and exploratory), not general; the
%       probe is better for support-style questions, not better in
%       general. The probe itself is not claimed as new.
%   C5  A matched-pair, multi-organisation spoken customer-support benchmark
%       in which retrieval need is not recoverable from surface form
%       (leak-test AUC 0.576 [0.541, 0.617]), released with the open-domain
%       speech set and all analysis code \PEND{release plan and URL}.
% NUMBERS: g^+ = 0.720, g^- = 0.592 [V1T retrieval_effect];
%   0.576 [0.541, 0.617] [branch v2/d2-support-a leak_test.json].
% DO NOT USE: "unified multi-timescale framework", "architectural strength"
%   (R1 minor), "Speculative RAG" (name of an unrelated method).
\todo{P-I.3: background stated as adopted
\cite{jiang2023flare,su2024dragin,arora2026streamrag,jung2026hiding,kadavath2022language},
then contributions C1--C5 as listed in the comment above, hedged where
novelty is claimed.}

% ---- P-I.4 ------------------------------------------------------------------
% PURPOSE: one-sentence roadmap with section references.
\todo{P-I.4: roadmap (Sec.~\ref{sec:related} to Sec.~\ref{sec:conclusion}).}
